\documentclass[11pt]{article}

\usepackage[margin=1in]{geometry}
\usepackage{amsmath,amssymb}
\usepackage{booktabs,longtable,tabularx}
\usepackage[round,authoryear]{natbib}
\usepackage{xcolor}
\usepackage[hidelinks]{hyperref}

\newcommand{\R}{\mathbb{R}}
\newcommand{\E}{\mathbb{E}}
\newcommand{\Var}{\operatorname{Var}}
\newcommand{\ESS}{\operatorname{ESS}}
\newcommand{\given}{\,\middle|\,}
\newcommand{\code}[1]{\texttt{#1}}
\newcommand{\cmark}{\textsf{yes}}
\newcommand{\pmark}{\textsf{partial}}
\newcommand{\xmark}{\textsf{no}}

\title{Hamiltonian Monte Carlo for Multimodal Targets:\\
Mechanisms, Corrections, Diagnostics, and a BayesFilter Evaluation Strategy}
\author{BayesFilter research note}
\date{9 August 2026}

\begin{document}
\maketitle

\begin{abstract}
Standard Hamiltonian Monte Carlo (HMC) can explore a connected typical set very
efficiently while failing almost completely to cross a low-density barrier between
isolated modes.  This failure can coexist with finite trajectories, high acceptance,
and apparently satisfactory within-basin convergence diagnostics.  This survey
separates three tasks that are often conflated: discovering important modes, moving
between known modes, and estimating their relative posterior masses.  It reviews
replica exchange, tempered transitions, symmetric-schedule tempered HMC,
continuous and geometric tempering, wormhole geometry, transport maps, annealed
importance sampling (AIS), and sequential Monte Carlo (SMC).  Exactness requires
the appropriate product target, Metropolis correction, Jacobian, or importance
weight; observing transitions is not itself an estimate of relative mode mass.  For
BayesFilter, the recommended first comparison is fixed-kernel parallel-tempered HMC
as an exact baseline and symmetric-schedule tempered HMC as a single-chain
candidate, with AIS or annealed SMC supplying independent evidence about global
mass.  NeuTra remains useful for within-region geometry but is not, by itself, a
multimodal coverage or weight authority.  The recommendation is an experiment
sequence, not a claim that any method is already correct for the SSL--LSTM target.
\end{abstract}

\section{Scope and source method}

The question is not whether multimodal sampling has a universal solution.  It is:
which HMC-compatible mechanisms can plausibly generate an empirical posterior with
the correct global mixture, what correction makes each method target the intended
distribution, and what evidence would reveal failure on the BayesFilter SSL--LSTM
problem?  This is a bounded technical survey, not a systematic review or a ranking
supported by a common benchmark.

The claim-bearing primary sources were stored locally, converted to text, and read
through their method, invariant-target or detailed-balance arguments, experiments,
and limitations.  The exact files, SHA-256 values, sections inspected, original-code
links, and one quarantined invalid download are recorded in
\code{.localresources/papers/multimodal\_hmc/CORPUS\_AUDIT.md}.  The original
1996 tempered-transitions article could not be recovered as a valid local PDF; the
algorithm and acceptance formula used below were therefore checked in the complete
technical account in Section~7 of \citet{neal2001ais}, which explicitly relates it to
\citet{neal1996tempered}.  No local implementation is claimed for any surveyed
method.

The BayesFilter policy in this work excludes TFP NUTS.  ``HMC'' below means a
fixed, explicitly tuned leapfrog kernel unless a cited method defines a different
corrected trajectory.  Literature evidence can nominate an implementation and an
experiment; it cannot establish correctness, speed, or posterior validity for this
repository.

\section{The failure mechanism in ordinary HMC}

Let the target have unnormalised density
\begin{equation}
  \widetilde\pi(q)=\exp\{-U(q)\}, \qquad q\in\R^d,
\end{equation}
and augment it with momentum $p\sim N(0,M)$.  Euclidean HMC uses
\begin{equation}
  H(q,p)=U(q)+K(p), \qquad K(p)=\tfrac12p^\mathsf{T}M^{-1}p.
\end{equation}
Exact Hamiltonian flow preserves $H$; a reversible, volume-preserving numerical
trajectory is corrected with
\begin{equation}
  a\{(q,p),(q',p')\}=1\wedge
  \exp\{-H(q',p')+H(q,p)\}.
  \label{eq:hmc-accept}
\end{equation}
Momentum refresh changes the energy between trajectories, but its kinetic energy is
typically of order $d/2$, with fluctuations of order $\sqrt d$.  A potential barrier
far into the kinetic-energy tail is therefore crossed only rarely
\citep{betancourt2017conceptual,graham2017continuous,nishimura2020geometric}.

This explains an important distinction.  HMC may mix rapidly \emph{inside} each
mode and still have a transition time between modes far longer than the run.  Low
rank-normalised $\widehat R$ cannot rule that out when every chain is confined to the
same basin, when initialisation omits an important basin, or when a symmetry is not
visible in the monitored marginal.  Conversely, deliberately putting equal numbers
of chains in two basins does not imply that the posterior masses are equal.  Their
initialisation counts were chosen by the analyst, not sampled from the posterior.

For later discussion, define three separate objectives:
\begin{description}
  \item[Discovery.] Locate every basin with non-negligible posterior mass.
  \item[Transition.] Move repeatedly among those basins while preserving a declared
        target.
  \item[Weight.] Estimate the relative posterior mass of the basins, with Monte Carlo
        uncertainty.
\end{description}
No finite generic algorithm certifies discovery of every important mode.  A method
that improves transition can still give poor weights if it has not reached its global
stationary regime.  A weighted path method can estimate masses consistently while
having intolerably variable weights in finite computation.

\section{Discrete tempering}

\subsection{Bridging family}

Let $g_0$ be a normalised base density and define inverse temperatures
$0=\beta_0<\cdots<\beta_L=1$ with
\begin{equation}
  \widetilde\pi_\beta(q)
  =g_0(q)^{1-\beta}\widetilde\pi(q)^\beta,
  \qquad
  \pi_\beta(q)=\frac{\widetilde\pi_\beta(q)}{Z_\beta}.
  \label{eq:bridge}
\end{equation}
For a pure power bridge, $g_0$ is omitted.  Smaller $\beta$ flattens target-energy
barriers, but it can also change mode volumes and therefore change relative mode
probabilities.  A broad mode becoming dominant at high temperature can make a
narrow mode effectively disappear.  This unequal-scale failure grows rapidly with
dimension and affects parallel tempering, tempered HMC, and tempered SMC
\citep{park2025athmc}.

\subsection{Parallel tempering or replica exchange}

Replica exchange runs one invariant Markov kernel at each temperature and targets
the exact product density
\begin{equation}
  \Pi(q_0,\ldots,q_L)=\prod_{\ell=0}^{L}\pi_{\beta_\ell}(q_\ell).
\end{equation}
An exchange of adjacent states $q_\ell$ and $q_{\ell+1}$ is accepted with
\begin{equation}
  a_{\mathrm{swap}}=1\wedge
  \frac{widetilde\pi_{\beta_\ell}(q_{\ell+1})
        \widetilde\pi_{\beta_{\ell+1}}(q_\ell)}
       {\widetilde\pi_{\beta_\ell}(q_\ell)
        \widetilde\pi_{\beta_{\ell+1}}(q_{\ell+1})}.
  \label{eq:swap}
\end{equation}
This is a detailed-balance correction for the product target, so samples at
$\beta_L=1$ target the original posterior.  \citet{hukushima1996exchange} show that
exchange probabilities fall exponentially as the temperature gap grows and identify
three distinct operating conditions: every adjacent exchange must remain possible,
replicas must traverse the full temperature range, and a replica at the hot end must
actually forget its original basin.  The first two do not imply the third.

For HMC, each within-temperature kernel can be an ordinary fixed HMC transition
with the temperature-specific potential and its own Metropolis correction.  This is
the simplest exact multimodal baseline for BayesFilter because it changes the
orchestration more than the local kernel.  Its costs are multiple concurrent states,
temperature-specific tuning, and a ladder that may require many replicas as dimension
increases.  Swap acceptance alone is explanatory, not proof of global mixing.

\subsection{Tempered transitions}

Tempered transitions propose one global move by applying an upward sequence of
kernels through increasingly diffuse distributions and a matched downward sequence
back to the target.  Suppose $\widehat T_j$ leaves $\pi_j$ invariant and
$\check T_j$ is its reversal with respect to $\pi_j$.  Starting at
$\widehat x_0$, the construction produces
\begin{equation*}
 \widehat x_0\rightarrow\widehat x_1\rightarrow\cdots
 \rightarrow\widehat x_{L-1}\rightarrow\bar x_L
 \rightarrow\check x_{L-1}\rightarrow\cdots\rightarrow\check x_0.
\end{equation*}
The terminal proposal is accepted with the telescoping density-ratio product given
by Eq.~(27) of \citet{neal2001ais}.  In compact notation,
\begin{equation}
 a_{\mathrm{TT}}=1\wedge
 \prod_{j=1}^{L}
 \frac{\widetilde\pi_j(\widehat x_{j-1})
       \widetilde\pi_{j-1}(\check x_{j-1})}
      {\widetilde\pi_{j-1}(\widehat x_{j-1})
       \widetilde\pi_j(\check x_{j-1})}.
 \label{eq:tt}
\end{equation}
The paired path and correction preserve the target; merely annealing upward and
back without Eq.~\eqref{eq:tt} does not.  Each proposal costs both an upward and a
downward traversal.  The ladder can still be deceptive: intermediate distributions
may suppress a target-relevant mode or fail to connect the regions that matter.

\section{Tempering inside Hamiltonian dynamics}

\subsection{Symmetric-schedule tempered HMC}

\citet{park2025athmc} replace the fixed Hamiltonian during a proposal by
\begin{equation}
 H_\alpha(q,p,t)=\tfrac12p^\mathsf{T}M^{-1}p+\frac{U(q)}{\alpha(t)},
 \label{eq:park-hamiltonian}
\end{equation}
where the temperature $\alpha(t)$ rises and then falls.  The outward part helps a
trajectory leave a basin; the return to unit temperature directs the endpoint back
toward a target mode.  The numerical step size is scaled using a transformed time
coordinate.  Crucially, the schedule is symmetric, starts and ends at unit
temperature, and the terminal state is accepted using the original endpoint
Hamiltonians as in Eq.~\eqref{eq:hmc-accept}.  Proposition~1 and its appendix proof
show that this symmetric construction is reversible and leaves $\pi$ invariant.

This method is attractive for BayesFilter because it produces a single fixed-length
HMC-style transition without a permanent population of replicas.  It is not made
exact by a visually symmetric temperature plot alone: half-step indexing, numerical
reversibility, endpoint temperature, and Metropolis correction are part of the
construction.  Its maximum temperature, schedule length, transformed step-size
scaling, and local mass matrix need target-specific tuning.  The paper explicitly
shows the unequal-mode-scale difficulty discussed after Eq.~\eqref{eq:bridge}.
Transitions therefore nominate a viable sampler but do not by themselves establish
relative masses.

\subsection{Continuously tempered HMC}

\citet{graham2017continuous} augment $q$ with a continuous temperature-control
coordinate $u$.  A smooth map $\beta(u)$ connects the target to a normalised base,
so that the extended potential contains
\begin{equation}
  \beta(u)U(q)+(1-\beta(u))[-\log g_0(q)]
\end{equation}
plus a bias or confinement term for $u$.  Ordinary leapfrog on the augmented
Hamiltonian remains reversible and volume preserving, followed by a Metropolis
correction.  Target expectations and the target normalising constant are recovered
with the paper's conditional or importance-weighted estimators, rather than by
silently treating all augmented states as target draws.

The method unifies temperature movement and state movement in one differentiable
dynamics.  Its practical dependence on a base distribution that overlaps the target,
and on calibration of the temperature marginal or normalising-constant estimate, is
substantial.  A base trained on one mode can provide an easy route to that mode and
no useful route to another.  The thermostat-assisted extension of
\citet{luo2019thermostat} adds Nos\'e--Hoover variables and adaptive biasing force for
minibatch noise.  That extension is scientifically relevant but is not the first
BayesFilter candidate: it changes the dynamics and stochastic-gradient assumptions
more than the fixed-HMC problem requires.

\subsection{Geometric tempering}

Riemannian HMC with a position-dependent metric $G(q)$ uses
\begin{equation}
  H(q,p)=-\log\pi(q)+\tfrac12\log|G(q)|
         +\tfrac12p^\mathsf{T}G(q)^{-1}p.
  \label{eq:rmhmc}
\end{equation}
The metric terms cancel appropriately after integrating over momentum, preserving
the intended position marginal.  \citet{nishimura2020geometric} choose metrics whose
volume factor satisfies
\begin{equation}
  |G_T(q)|^{1/2}\ \propto\ \pi(q)^{1-1/T},
  \label{eq:gthmc-condition}
\end{equation}
which lowers effective energy barriers by a factor $T$.  Their isometric variant
changes distances in all directions; the directional variant can be more efficient
when a useful inter-mode direction is known.  Barrier reduction is not sufficient to
cause a transition, however, and the directional construction imports geometric
knowledge.

Fast changes in the metric make standard integrators unstable.  The paper therefore
develops a reversible variable-step integrator and includes its non-unit Jacobian in
the acceptance probability.  This is not a drop-in mass-matrix change.  A faithful
BayesFilter implementation would need the metric determinant, its derivatives, the
integrator, Jacobian, and reversibility tests.  The larger implementation and tuning
surface makes GTHMC a later candidate, not the initial baseline.

\section{Mode-informed geometry and transport}

\subsection{Wormhole HMC}

Wormhole HMC shortens the metric distance along tubes connecting known mode
locations \citep{lan2014wormhole}.  The method adds a vector field, an auxiliary
dimension to reduce interference among connections, and a correction for a
non-volume-preserving integrator.  Its appendices establish detailed balance and
stationarity for the corrected transition.  The paper begins with known or
approximately known modes, then proposes residual-density optimisation to find new
ones.  Adaptation of the mode library and wormhole network occurs at regeneration
times so that arbitrary ongoing adaptation does not destroy stationarity.

This construction directly encodes useful global information, which can make it
powerful when reliable mode locations are available.  It is also the source of its
main limitation: missed modes, inaccurate local geometry, a poorly connected
wormhole network, or rare regenerations can defeat the method.  The approximate
weights used in its residual search are search machinery, not an independent
posterior-mass authority.  It is best classified as a mode-discovery and
mode-transition mechanism, followed by separate weight validation.

\subsection{Transport maps and NeuTra}

For a bijection $q=f(z)$, exact change of variables gives
\begin{equation}
  \pi_Z(z)=\pi(f(z))\left|\det Df(z)\right|.
  \label{eq:change-variable}
\end{equation}
HMC or another Metropolis-corrected kernel targeting Eq.~\eqref{eq:change-variable}
is exact even when $f$ only approximately Gaussianises the target.  Transport-map
MCMC combines learned maps with Metropolis--Hastings correction and gives conditions
under which adaptive, inexact maps retain the exact target
\citep{parno2018transport}.  NeuTra learns an inverse-autoregressive-flow map by
minimising the reverse divergence $\mathrm{KL}(q\Vert\pi)$, runs HMC in $z$-space,
then maps draws forward \citep{hoffman2019neutra}.

The correction fixes proposal bias conditional on the chain exploring the transformed
target; it does not make a missing mode appear in the training distribution.  Reverse
KL can assign almost no proposal mass to a separated mode with little training signal
from that mode.  A chain can then have excellent local geometry and still have a
global transition problem.  Consequently, NeuTra should be used as within-region
preconditioning, or trained only after an independent global method has supplied
weighted multimodal coverage.  NeuTra loss, validation loss, and local HMC
acceptance are explanatory diagnostics, not multimodal posterior-weight evidence.

\section{Estimating global mass: AIS and SMC}

Mode transition counts estimate posterior masses only after global stationarity and
adequate effective sample size have been established.  When that premise is exactly
what is in doubt, a separate weighted path provides discriminating evidence.

\subsection{Annealed importance sampling}

Let $f_0,\ldots,f_L$ be unnormalised densities with normalising constants
$Z_0,\ldots,Z_L$.  AIS starts at $x_0\sim\pi_0$ and applies transitions that leave
successive $\pi_j$ invariant.  With a consistent indexing convention, its weight is
\begin{equation}
  W=\prod_{j=1}^{L}\frac{f_j(x_{j-1})}{f_{j-1}(x_{j-1})}.
  \label{eq:ais-weight}
\end{equation}
The average of independent AIS weights is an unbiased estimator of $Z_L/Z_0$;
self-normalised weights estimate target expectations \citep{neal2001ais}.  Terminal
mode counts need not be representative because the weights, not unweighted counts,
perform the correction.  This makes AIS useful for estimating relative basin masses:
combine the terminal basin indicator with $W$ and report uncertainty across
independent runs.

The guarantee does not imply useful finite variance.  If an important mode is almost
never found, the few runs that reach it can carry enormous weights, and the upward
tail may be absent from the observed sample.  Neal explicitly warns of heavy-tailed
importance weights that yield drastically wrong finite results without an obvious
warning.  Weight ESS, maximum normalised weight, repeated independent batches, and
schedule sensitivity are therefore mandatory diagnostics.  No observed mode at all
is a discovery failure, not evidence of zero posterior mass.

\subsection{Annealed sequential Monte Carlo}

SMC samplers propagate a population through a sequence such as
Eq.~\eqref{eq:bridge}, alternating importance weighting, unbiased resampling, and
MCMC rejuvenation kernels invariant for the current target
\citep{delmoral2002smc,delmoral2006smc}.  Resampling prevents all future computation
from being spent on negligible-weight particles; rejuvenation restores diversity
after replication.  Products of mean incremental weights estimate normalising
constants.

The particle population can preserve multiple basins better than one local chain,
but only if the schedule resolves changes in relative mass and rejuvenation moves
particles enough within relevant regions.  Premature resampling can delete a basin;
later invariant kernels cannot recover it if inter-mode movement is absent.  Adaptive
temperature placement based on conditional ESS can control a particular weight
degeneracy measure but does not certify mode discovery.  Results must include basin
survival across stages, ancestry diversity, weight ESS, rejuvenation movement, and
replicate uncertainty.  Unequal mode scales remain a difficult case
\citep{park2025athmc}.

\section{What each family actually addresses}

Table~\ref{tab:comparison} is a functional classification, not a performance ranking.
``Exact target'' means that the cited correction is present and its assumptions hold;
it does not mean a finite run has mixed.

\small
\begingroup
\setlength{\tabcolsep}{3pt}
\begin{longtable}{p{0.18\linewidth}p{0.13\linewidth}p{0.12\linewidth}p{0.12\linewidth}p{0.16\linewidth}p{0.20\linewidth}}
\caption{Roles and principal failure modes.}\label{tab:comparison}\\
\toprule
Family & Exact target mechanism & Discovery & Transition & Relative weight & Principal finite-run risk \\
\midrule
\endfirsthead
\toprule
Family & Exact target mechanism & Discovery & Transition & Relative weight & Principal finite-run risk \\
\midrule
\endhead
Ordinary fixed HMC & Endpoint Metropolis correction & \xmark & local & only after global mixing & Energy barrier never crossed \\
Parallel tempering & Product target plus swap correction & \pmark & \cmark & target-replica occupancy after round trips & Ladder bottleneck; hot chain retains basin identity \\
Tempered transitions & Reversed up/down kernels plus density-ratio correction & \pmark & \cmark & only after global mixing; optional AIS estimators & Deceptive bridge; low global acceptance \\
Symmetric THMC & Symmetric schedule and endpoint Metropolis correction & \pmark & \cmark & only after global mixing & Wrong schedule discretisation; unequal mode scales \\
Continuous tempering & Corrected augmented target and target-specific weighting and conditioning & \pmark & \cmark & possible through weighted estimators & Poor base overlap; temperature-marginal calibration \\
Geometric tempering & RMHMC determinant, reversible integrator, and Jacobian correction & \xmark/\pmark & \cmark & only after global mixing & Metric/integrator complexity; missing direction \\
Wormhole HMC & Corrected non-volume-preserving dynamics; regeneration-safe adaptation & \pmark & \cmark & only after global mixing & Requires or must discover modes and connections \\
Transport/NeuTra HMC & Exact transformed density and HMC/MH correction & \xmark & \pmark & only after global mixing & Learned map omits a mode \\
AIS & Importance weights across invariant bridges & \pmark & not a chain property & \cmark & Rare enormous weights; important mode never found \\
Annealed SMC & Sequential weights, unbiased resampling, invariant rejuvenation & \pmark & \pmark & \cmark & Basin extinction and particle degeneracy \\
\bottomrule
\end{longtable}
\endgroup
\normalsize

\section{Diagnostics that answer different questions}

Diagnostics must be assigned a role before looking at results.

\paragraph{Hard validity vetoes.}
Any non-finite target, state, gradient, correction term, or weight; an invalid target
status; an uncorrected non-volume-preserving map; failure of a reversibility or
invariance fixture; a divergence under the declared energy-error rule; or a corrupted
artifact invalidates that run.  Acceptance is not a convergence gate, although an
extreme value can trigger tuning investigation.

\paragraph{Mode-discovery diagnostics.}
On synthetic targets with known modes, define basins before the claim run and report
whether every material basin was reached from dispersed initial states.  On unknown
targets, use independent optimisation, hot replicas, AIS/SMC populations, or multiple
discovery seeds, but label the resulting list as incomplete.  Finding the same list
twice is descriptive stability, not proof of completeness.

\paragraph{Transition diagnostics.}
For parallel tempering, report adjacent swap rates, temperature-index traces,
complete cold-to-hot-to-cold round trips, and whether configurations lose their
starting-basin label at the hottest temperature.  For any sampler, report basin
transition counts, transition directions, dwell times, and results from each
initial basin.  A few transitions show reachability; they do not establish stationary
weights.

\paragraph{Weight diagnostics.}
For target-replica occupancy, give binomial or Markov-chain-aware uncertainty using
effective rather than raw counts.  For AIS/SMC, report weight ESS,
\begin{equation}
  \ESS_W=\frac{(\sum_i W_i)^2}{\sum_i W_i^2},
\end{equation}
the largest normalised weight, basin-specific weighted mass with uncertainty, and
replication across independent batches.  Check schedule refinement and alternative
bases.  Agreement of two methods with the same undiscovered-mode failure is not an
independent authority.

\paragraph{Within-mode diagnostics.}
Modern split/folded rank $\widehat R$, bulk/tail ESS, energy error, movement, and
posterior geometry checks remain necessary.  They are conditional on the explored
support and cannot be promoted to a global coverage claim.

\paragraph{Downstream predictive diagnostics.}
For synthetic SSL--LSTM data, the relevant object is the posterior predictive
mixture
\begin{equation}
 p(y\mid y_{\mathrm{obs}})
 =\int p(y\mid\theta)p(\theta\mid y_{\mathrm{obs}})\,d\theta.
 \label{eq:posterior-predictive}
\end{equation}
Its empirical simulation must independently select one retained posterior draw for
each output path and generate one conditional path from that selected row.  Reusing
one posterior mean or one randomly chosen parameter for many paths computes a
different distribution.  Predictive agreement is scientifically useful when the
parameterisation is non-identifiable, but it does not prove posterior correctness:
different parameter distributions can induce the same finite-horizon output law.

\section{BayesFilter staged recommendation}

The immediate target is a valid empirical posterior mixture for the SSL--LSTM
posterior-predictive diagnostic, not the maximum possible number of mode switches.
The following sequence keeps local geometry, global transitions, and relative weights
separate.

\begin{enumerate}
  \item \textbf{Build known-answer fixtures.}  Before the SSL--LSTM target, use
  unequal-weight Gaussian mixtures with equal and unequal component scales.  Check
  target density, gradients, XLA parity, reversibility, swap or endpoint corrections,
  basin probabilities, and uncertainty calibration.  A sampler that cannot recover
  known unequal weights is not eligible for the real target.

  \item \textbf{Use parallel-tempered fixed HMC as the exact baseline.}  Keep the
  current fixed HMC kernel at every temperature; tune the ladder and local kernel on
  calibration data or calibration seeds, freeze them, and run untouched validation.
  Require adjacent communication, repeated temperature round trips, hot-basin
  forgetting, finite corrected transitions, and target-replica convergence from each
  known basin.  Do not infer weights from equal chain initialisation.

  \item \textbf{Evaluate symmetric-schedule THMC as the plain candidate.}  Implement
  the exact symmetric half-step schedule and endpoint Metropolis correction from
  \citet{park2025athmc}.  Compare it with parallel tempering at matched target/gradient
  evaluation budgets.  Predeclare the primary comparison and uncertainty analysis;
  transition count and runtime alone are descriptive.

  \item \textbf{Add an independent weighted path.}  Run AIS first because it is the
  smallest separate normalising-constant and basin-mass authority.  If weights
  collapse, move to annealed SMC with conditional-ESS temperature placement,
  resampling, and fixed-HMC rejuvenation.  Keep independent runs and report
  basin-specific weighted uncertainty.  A low weight ESS is a repair trigger, not a
  reason to use unweighted terminal counts.

  \item \textbf{Use NeuTra only where it has evidence.}  A trained transport may
  accelerate each local or tempered HMC kernel.  It is not allowed to create a
  multimodal-weight authority from reverse-KL training loss or a mode-limited replay
  set.  If a transport is trained from globally weighted particles, test whether it
  preserves both basin coverage and corrected HMC movement.

  \item \textbf{Issue the posterior archive only after agreement.}  The archive must
  exclude warm-up, bind the exact target and sampler result, record equal or explicit
  importance-weight semantics, and classify relative mode weights as resolved only
  when the predeclared global evidence passes.  If transition and weighted evidence
  conflict, preserve both artifacts and repair; do not average them into a claim.

  \item \textbf{Then run the posterior-predictive diagnostic.}  Draw one posterior
  parameter independently per path as in Eq.~\eqref{eq:posterior-predictive}, simulate
  the truth arm independently, and compare complete-path distributions at each
  horizon.  Non-rejection means ``not distinguished at this sample size and horizon,''
  not equality and not HMC correctness.
\end{enumerate}

The initial baseline ladder is methodologically justified, but its numerical ladder,
temperature ceiling, step sizes, leapfrog counts, number of replicas, AIS schedule,
and compute budget are deliberately not specified here.  Those are target-specific
hypotheses requiring a reviewed experiment plan and calibration canaries; literature
defaults cannot be silently promoted to the SSL--LSTM target.

\section{Decision and nonclaims}

\begin{center}
\begin{tabularx}{\textwidth}{p{0.22\textwidth}X}
\toprule
Decision field & Status \\
\midrule
Primary conclusion & Ordinary HMC evidence is insufficient for a posterior known to have an unvisited competing mode. \\
Viable first baseline & Parallel-tempered fixed HMC with exact swap correction and global-mixing diagnostics. \\
Viable plain candidate & Symmetric-schedule, endpoint-corrected THMC. \\
Required independent evidence & AIS or annealed SMC basin-mass estimates with weight diagnostics and replicate uncertainty. \\
NeuTra status & Within-region geometry aid; not a standalone discovery or mode-weight authority. \\
Main uncertainty & Whether all material SSL--LSTM modes can be found and whether the annealing path avoids unequal-scale weight collapse. \\
Next justified action & Write and audit a known-mixture plus SSL--LSTM tempering experiment plan before implementation or a long run. \\
Not concluded & No surveyed method is promoted, no SSL--LSTM posterior is validated, and no runtime or statistical ranking is established. \\
\bottomrule
\end{tabularx}
\end{center}

The strongest alternative explanation for the current apparent multimodality is an
implementation or coordinate error rather than a genuine second posterior region.
That possibility must be checked by evaluating both modes with the same canonical
target and gradients before a sampler comparison.  Conversely, successful transitions
between two known regions would not rule out a third undiscovered region.  The weakest
part of all finite multimodal evidence is mode completeness; the proposed weighted
methods address relative mass among reached regions, not a proof that no other region
exists.

\bibliographystyle{plainnat}
\bibliography{multimodal_hmc_survey}

\end{document}
